\documentclass[../DoAn.tex]{subfiles}
\begin{document}

\section{Đặt vấn đề}
\label{section:1.1}

Vài năm gần đây, tài chính phi tập trung (Decentralized Finance - DeFi)~\cite{werner2022sok} phát triển nhanh chóng, và một trong những thành phần quan trọng nhất của nó là các sàn giao dịch phi tập trung (Decentralized Exchange - DEX). Khác với sàn giao dịch tập trung, DEX cho phép người dùng đổi tài sản số trực tiếp với nhau thông qua hợp đồng thông minh, mà không phải gửi tài sản cho một bên trung gian giữ hộ. Hiện nay khối lượng giao dịch mỗi ngày trên các DEX lớn như Uniswap đã lên tới hàng tỉ đô la, cho thấy mô hình này ngày càng được sử dụng nhiều.

Nếu nhìn lại quá trình phát triển của DEX, có thể thấy nó được cải tiến dần qua từng giai đoạn. Những DEX đầu tiên, tiêu biểu là Uniswap, hoạt động theo cơ chế tạo lập thị trường tự động (Automated Market Maker - AMM). Với cơ chế này, người dùng giao dịch trực tiếp với một hồ thanh khoản (liquidity pool), giá được tính theo công thức dựa trên lượng tài sản trong hồ. Mô hình này đơn giản và luôn sẵn sàng giao dịch, nhưng có một điểm yếu: giá người dùng nhận được phụ thuộc hoàn toàn vào lượng thanh khoản hiện có trong hồ trên chuỗi. Khi giao dịch một khối lượng lớn, giá sẽ bị lệch đáng kể so với mong đợi (hiện tượng trượt giá), và điều này còn tạo cơ hội cho một số bên lợi dụng để kiếm lời từ chính giao dịch của người dùng (hiện tượng MEV~\cite{daian2020flashboys}, trình bày thêm ở Phụ lục~\ref{apdx:mev}).

Để khắc phục, gần đây xuất hiện một hướng tiếp cận mới gọi là giao dịch dựa trên ý định (intent-based). Ở hướng này, người dùng không tự thực hiện giao dịch mà ký một ``ý định'' mô tả tài sản muốn đổi. Việc tìm nguồn thanh khoản và thực hiện giao dịch do các bên thứ ba (gọi là filler) cạnh tranh nhau đảm nhận. Ưu điểm của hướng này là filler không bị giới hạn trong một hồ thanh khoản nào; họ có thể tập hợp thanh khoản từ nhiều nguồn khác nhau, kể cả các nguồn ngoài chuỗi (off-chain), nên người dùng có cơ hội nhận được giá tốt hơn so với khi giao dịch trong phạm vi một hồ AMM.

Tuy nhiên, khi tìm hiểu các hệ thống đang có (trình bày kỹ hơn ở Chương~\ref{chapter:Related_works}), có thể thấy ba vấn đề chưa được giải quyết. Thứ nhất, hầu hết các hệ thống khớp lệnh theo nguyên tắc toàn bộ hoặc không có gì (all-or-nothing), tức là khớp trọn vẹn cả lệnh hoặc không khớp gì, mà chưa hỗ trợ khớp từng phần. Thứ hai, filler đăng ký nhận lệnh nhưng nếu sau đó không thực hiện thì cũng không chịu chế tài nào, nên người dùng có thể phải chờ lâu hoặc lệnh bị bỏ dở cho tới khi hết hạn. Thứ ba, việc giao dịch giữa hai chuỗi khác nhau (xuyên chuỗi) thường phải đi qua cầu nối (bridge), mà cầu nối là nơi giữ tài sản của người dùng nên dễ trở thành mục tiêu tấn công.

Từ đó, đồ án đặt mục tiêu xây dựng một hệ thống khớp lệnh từng phần. Hệ thống buộc filler đặt cọc và chịu phạt nếu bỏ dở, đảm bảo lệnh luôn được hoàn tất trước hạn, đồng thời hỗ trợ giao dịch xuyên chuỗi không cần cầu nối.

\section{Mục tiêu, phạm vi và giới hạn}
\label{section:1.2}

\subsection{Mục tiêu}
\label{subsection:1.2.1}

Đồ án đặt ra năm mục tiêu, tất cả đều nằm ở tầng hợp đồng thông minh - phần lõi của hệ thống:

\begin{itemize}
    \item Xây dựng cơ chế khớp lệnh \textbf{từng phần} dựa trên \textbf{đấu giá kiểu Hà Lan}, trong đó giá giảm dần theo từng khối để hạn chế bị khai thác MEV.
    \item Cho phép người dùng ký lệnh ngoài chuỗi bằng chuẩn \textbf{EIP-712} kết hợp \textbf{Permit2}, nhờ đó không tốn phí gas khi tạo lệnh.
    \item Xây dựng cơ chế \textbf{đặt cọc và phạt} để ràng buộc trách nhiệm của filler khi nhận lệnh.
    \item Xây dựng một \textbf{cơ chế dự phòng} để chắc chắn phần lệnh còn lại luôn được khớp trước khi hết hạn.
    \item Hỗ trợ \textbf{giao dịch xuyên chuỗi} giữa hai mạng tương thích EVM mà không dùng cầu nối, đồng thời vẫn cho phép khớp từng phần.
\end{itemize}

\subsection{Phạm vi}
\label{subsection:1.2.2}

Đồ án tập trung chủ yếu vào tầng \textbf{hợp đồng thông minh}, vì đây là nơi chứa toàn bộ logic quan trọng của hệ thống. Phần backend và frontend được xây dựng ở mức đủ để chạy thử và minh họa, không hướng tới một sản phẩm hoàn chỉnh. Tài sản giao dịch được giới hạn ở các token theo chuẩn ERC-20. Hệ thống được triển khai và kiểm thử trên mạng giả lập cục bộ tạo bằng Anvil: một mạng sao chép trạng thái từ Ethereum mainnet cho trường hợp giao dịch trong cùng một chuỗi, và hai mạng giả lập riêng cho trường hợp xuyên chuỗi.

Phần chiến lược tính lợi nhuận của filler không được xử lý trong đồ án. Trong thực tế, filler có thể là các nhà giao dịch, nhóm nghiên cứu hoặc công ty hoạt động độc lập, và họ thường không công khai mã nguồn cũng như cách quyết định có tham gia khớp lệnh hay không. Vì vậy, đồ án tập trung xây dựng các quy tắc trên hợp đồng (đặt cọc, phạt, hoàn cọc) để ràng buộc filler, còn việc filler tối ưu lợi nhuận nằm ngoài phạm vi nghiên cứu. Hai filler kèm theo trong đồ án là các chương trình mô phỏng đơn giản dùng để chạy thử các luồng giao dịch, không phải filler thực tế.

\subsection{Giới hạn}
\label{subsection:1.2.3}

Đồ án hiện được chạy và kiểm thử trên mạng giả lập, chưa triển khai lên mạng chính thức (mainnet). Để vận hành trên thực tế, hệ thống cần thêm nhiều điều kiện mà đồ án chưa có: thanh khoản thực, filler thực, một đợt kiểm toán bảo mật do chuyên gia thực hiện, cùng chi phí gas tương ứng. Ngoài ra, một số phần trong thiết kế vẫn mang tính tập trung, ví dụ khóa đồng ký (cosigner) hiện do một bên duy nhất giữ. Các giới hạn này được trình bày chi tiết hơn ở Chương~\ref{chapter:conclusion}, kèm theo hướng khắc phục.

\section{Định hướng giải pháp}
\label{section:1.3}

Đồ án chọn hướng giao dịch dựa trên ý định: người dùng (swapper) ký một lệnh mô tả muốn đổi gì, còn việc tìm thanh khoản và thực hiện được giao cho các filler cạnh tranh nhau đảm nhận. Cách này tách ``ý muốn'' khỏi ``thực hiện'', giúp giảm chi phí tạo lệnh và cải thiện kết quả nhờ cạnh tranh.

Dựa trên hướng đó, đồ án xây dựng một hệ thống đặt tên là \textbf{NeutronX}, gồm bốn thành phần làm việc cùng nhau: tầng hợp đồng thông minh đảm nhận việc kiểm tra và thực hiện lệnh trên chuỗi; tầng backend đảm nhận việc đồng ký lệnh và theo dõi trạng thái; tầng filler đóng vai các bên khớp lệnh; và giao diện người dùng để tạo và theo dõi lệnh. Kiến trúc chi tiết của hệ thống được trình bày ở mục~\ref{sec:arch}.

Đóng góp chính gồm bốn cơ chế hợp đồng: khớp lệnh từng phần theo đấu giá Hà Lan để hạn chế MEV; đặt cọc và phạt để ràng buộc filler; cơ chế dự phòng đảm bảo lệnh luôn hoàn tất; và giao dịch xuyên chuỗi không cầu nối với khớp từng phần. Đây là sản phẩm thử nghiệm, kiểm chứng trên môi trường giả lập, không nhằm cạnh tranh với các hệ thống thương mại đang vận hành.

\section{Bố cục đồ án}
\label{section:1.4}

Phần còn lại được tổ chức như sau. Chương~\ref{chapter:Related_works} khảo sát DEX qua Uniswap, chỉ ra chỗ còn thiếu và trình bày ý tưởng NeutronX. Chương~\ref{chapter:Experiment} trình bày thiết kế và cài đặt chi tiết tầng hợp đồng, kèm tổng quan backend và frontend. Chương~\ref{chapter:Methodology} trình bày kiểm thử và rà soát bảo mật. Chương~\ref{chapter:conclusion} tổng kết kết quả, nêu hạn chế và đề xuất hướng phát triển.

\end{document}
